\section{Zadanie 8 (Oskar Prusik)}
Łamiemy kij długości \(1\) w jednostajnie losowo wybranym punkcie.
Niech \(X\) oznacza punkt złamania.


(i) Jaka jest oczekiwana długość krótszej części?  

(ii) Jaki jest oczekiwany stosunek długości części krótszej do części dłuższej?  

(iii) Jaki jest oczekiwany stosunek długości części dłuższej do części krótszej?


\subsection*{Rozwiązanie}
 Ponieważ wybór jest jednostajny na odcinku \([0,1]\), zmienna losowa \(X\) ma rozkład jednostajny na tym przedziale. Gęstość prawdopodobieństwa dana jest wzorem:

\[
f_X(x) = 
\begin{cases}
	1 & \text{dla } x \in [0,1] \\
	0 & \text{w pozostałych przypadkach}
\end{cases}
\]

Wprowadzamy oznaczenia:
\[
L = \min(X, 1-X), \quad D = \max(X, 1-X)
\]
gdzie \(L\) to długość krótszej części, a \(D\) to długość dłuższej części.

\subsection*{(i)}
\[
\mathbb{E}[L] = \int_{0}^{1} \min(x, 1-x) \cdot f_X(x) \, dx = \int_{0}^{1} \min(x, 1-x) \, dx
\]
Punkt \(x = 0.5\) jest punktem podziału:
\[
\mathbb{E}[L] = \int_{0}^{0.5} x \, dx + \int_{0.5}^{1} (1-x) \, dx
\]
Obliczamy:
\[
\int_{0}^{0.5} x \, dx = \left[ \frac{x^2}{2} \right]_{0}^{0.5} = \frac{0.25}{2} = 0.125
\]
\[
\int_{0.5}^{1} (1-x) \, dx = \left[ x - \frac{x^2}{2} \right]_{0.5}^{1} = (1-0.5) - (0.5 - 0.125) = 0.5 - 0.375 = 0.125
\]
\[
\mathbb{E}[L] = 0.125 + 0.125 = 0.25 = \frac14
\]

\subsection*{(ii)}
\[
R_1 = \frac{L}{D} = \frac{\min(x, 1-x)}{\max(x, 1-x)}
\]
\[
\mathbb{E}[R_1] = \int_{0}^{1} \frac{\min(x, 1-x)}{\max(x, 1-x)} \cdot f_X(x) \, dx = \int_{0}^{1} \frac{\min(x, 1-x)}{\max(x, 1-x)} \, dx
\]
Całkujemy po dwóch przedziałach:
\[
\mathbb{E}[R_1] = \int_{0}^{0.5} \frac{x}{1-x} \, dx + \int_{0.5}^{1} \frac{1-x}{x} \, dx
\]

Pierwsza całka:
\[
\int_{0}^{0.5} \frac{x}{1-x} \, dx = \int_{0}^{0.5} \left( \frac{1}{1-x} - 1 \right) dx
\]
\[
= \left[ -\ln(1-x) - x \right]_{0}^{0.5} = \left( -\ln 0.5 - 0.5 \right) - \left( -\ln 1 - 0 \right) = \ln 2 - 0.5
\]

Druga całka:
\[
\int_{0.5}^{1} \frac{1-x}{x} \, dx = \int_{0.5}^{1} \left( \frac{1}{x} - 1 \right) dx
\]
\[
= \left[ \ln x - x \right]_{0.5}^{1} = (\ln 1 - 1) - (\ln 0.5 - 0.5) = (-1) - (-\ln 2 - 0.5) = \ln 2 - 0.5
\]

Sumując:
\[
\mathbb{E}[R_1] = (\ln 2 - 0.5) + (\ln 2 - 0.5) = 2\ln 2 - 1
\]

\subsection*{(iii)}
\[
R_2 = \frac{D}{L} = \frac{\max(x, 1-x)}{\min(x, 1-x)}
\]
\[
\mathbb{E}[R_2] = \int_{0}^{1} \frac{\max(x, 1-x)}{\min(x, 1-x)} \cdot f_X(x) \, dx = \int_{0}^{1} \frac{\max(x, 1-x)}{\min(x, 1-x)} \, dx
\]
\[
\mathbb{E}[R_2] = \int_{0}^{0.5} \frac{1-x}{x} \, dx + \int_{0.5}^{1} \frac{x}{1-x} \, dx
\]

\noindent\textbf{Obliczenie pierwszej całki:}
\[
\int_{0}^{0.5} \frac{1-x}{x} \, dx = \int_{0}^{0.5} \left( \frac{1}{x} - 1 \right) dx
\]
Całka \(\int_{0}^{0.5} \frac{1}{x} \, dx\) jest rozbieżna, co pokazujemy przez granicę:
\[
\lim_{\varepsilon \to 0^+} \int_{\varepsilon}^{0.5} \frac{1}{x} \, dx = \lim_{\varepsilon \to 0^+} \left[ \ln x \right]_{\varepsilon}^{0.5} = \ln 0.5 - \lim_{\varepsilon \to 0^+} \ln \varepsilon = +\infty
\]
Nawet po odjęciu skończonej całki \(\int_{0}^{0.5} 1 \, dx = 0.5\), pierwsza całka dąży do nieskończoności.

\noindent\textbf{Obliczenie drugiej całki:}
\[
\int_{0.5}^{1} \frac{x}{1-x} \, dx = \int_{0.5}^{1} \left( \frac{1}{1-x} - 1 \right) dx
\]
\[
= \left[ -\ln(1-x) - x \right]_{0.5}^{1}
\]
W \(x=1\): \(-\ln(0) - 1\) - tutaj \(\ln 0 = -\infty\), ale liczymy granicę:
\[
\lim_{t \to 1^-} \left[ -\ln(1-t) - t \right]_{0.5}^{t} = \lim_{t \to 1^-} \left[ (-\ln(1-t) - t) - (-\ln 0.5 - 0.5) \right]
\]
\[
= \lim_{t \to 1^-} \left[ -\ln(1-t) - t + \ln 0.5 + 0.5 \right]
\]
Człon \(-\ln(1-t) \to +\infty\), więc całka od \(0.5\) do \(1\) też jest rozbieżna!

Zatem obie całki w \(\mathbb{E}[R_2]\) są rozbieżne do \(+\infty\), co potwierdza:
\[
\mathbb{E}[R_2] = \infty
\]

\subsection*{Odpowiedź}
\[
\boxed{\text{(i) } \frac14} \quad \boxed{\text{(ii) } 2\ln 2 - 1} \quad \boxed{\text{(iii) } \infty}
\]
